\section{Finite residue actions from one K\"ahler class}\label{sec:scalars}

The restriction criterion asks for equality under links on lower strata.
To build enough sections with all these equalities, we need finiteness of
the induced actions on each fixed section space.  On projective strata this
is the usual log pluricanonical representation theorem.  On nonprojective
surfaces it holds here because the links all come from the boundary of one
compact K\"ahler fourfold and therefore compare restrictions of one
ambient K\"ahler class.

Fix the common even degree \(q\) and lines \(L_Z\) from
\eqref{eq:stratum-semiample-map}.  For \(d=0,1,2\), form a groupoid
\(\mathcal G_d\) whose objects are the finitely many \(d\)-dimensional
strata.  In dimension zero use all identifications of points, with the
canonical zero-form generator \(1\).  In dimension one use all crepant
residue comparisons between the indicated curve pairs.  In dimension two
use only the groupoid generated by the links of three-dimensional strata
from Proposition~\ref{prop:restriction-link}, together with their inverses.
Here a comparison includes equality of the pulled-back invertible
meromorphic adjoint subsheaves, as specified there.

Such a comparison induces an isomorphism
\[
 H^0(Z_1,L_{Z_1}^k)\longrightarrow H^0(Z_2,L_{Z_2}^k)
\]
for every \(k>0\).  Pull a section to a resolution of the graph, use the
equality of invertible meromorphic subsheaves, and descend to the other
normal stratum by projection formula.  This transport is compatible with
composition, products of sections, and residue restriction along boundary
divisors whose centers are birational on both strata.  Lemma
\ref{lem:transport-lower-restrictions} proves the compatibility also when
a floor curve is contracted, before that case enters the construction.
The same
argument shows that evaluation after a graph pullback is evaluation in
the identified actual line fibers, a fact needed for generation later.

\begin{proposition}[Finite isotropy images]\label{prop:finite-isotropy}
For every \(d\in\{0,1,2\}\), every \(d\)-stratum \(Z\), and every fixed
integer \(k>0\), the image of the self-comparisons in \(\mathcal G_d\)
on \(H^0(Z,L_Z^k)\) is a finite group.
\end{proposition}

The restriction on \(\mathcal G_2\) is essential.  A nonprojective K3
surface with trivial canonical line can have an automorphism acting by
a non-root-of-unity scalar on its holomorphic two-form
\cite[Theorems~3.5 and~4.1]{McMullen2002}.  Thus the proposition would be false
for arbitrary bimeromorphic self-comparisons of nonprojective surfaces.

\subsection{Points, curves, and projective surfaces}

On a point the zero-form generator is \(1\) and every comparison acts
as the identity.  A normal compact curve is smooth, and a bimeromorphic
map is an isomorphism preserving the boundary coefficients.  For genus
at least two the automorphism group is finite.  For genus one, the
stabilizer of a nonempty finite boundary support is finite; with empty
boundary translations act trivially on holomorphic forms and the linear
automorphism group is finite.  For genus zero, nefness of the adjoint and
coefficients at most one require at least two marked points.  If there
are exactly two, both coefficients are one and the log line is trivial
with generator \((dt/t)^{\otimes qk}\).  Scaling fixes it, and exchange
has sign one in the even degree.  With at least three marked points the
stabilizer is finite.  This proves Proposition~\ref{prop:finite-isotropy}
in dimensions zero and one.

Suppose a surface stratum \(Z\) is Moishezon.  Its klt perturbation in
Proposition~\ref{prop:strata-adjunction} gives rational singularities.
Namikawa's projectivity criterion \cite[Corollary~6']{Namikawa2002}
makes the compact K\"ahler \(Z\) projective.  Its effective full boundary
is dlt and its adjoint is semiample.  The analytic comparison graphs are
algebraic by Chow, and their meromorphic equality is the usual crepant
B-birational equality for compatible canonical identifications.
Fujino--Gongyo's log pluricanonical representation theorem
\cite[Theorem~1.1]{FujinoGongyo2014} gives finite image in every fixed
Cartier degree.  It applies exactly in this projective case.

\subsection{The class attached to a nonprojective surface}

Let \(Z\) now be a non-Moishezon surface stratum, and let \(P\) be the
minimal smooth surface of a compact K\"ahler resolution of \(Z\).
Point blowdowns preserve K\"ahlerness, for example by the even-\(b_1\)
criterion \cite[Theorem~11]{Buchdahl1999}.  Algebraic dimension is
bimeromorphically invariant, so \(P\) is nonprojective and has algebraic
dimension zero or one.  It has a nonzero holomorphic two-form: otherwise
\(H^2(P,\R)=H^{1,1}(P,\R)\), and rational approximation of a K\"ahler
class followed by Kodaira embedding would make \(P\) projective.  Thus
\(\kappa(P)\geq0\).  The minimal model is neither rational nor ruled;
its bimeromorphic maps are automorphisms, and bimeromorphic maps between
such minimal surfaces are isomorphisms
\cite[Proposition~3.5]{ProkhorovShramov2017}.  Under a common resolution,
\(H^0(Z,L_Z^k)\) embeds as a finite-dimensional space of meromorphic
\(qk\)-pluricanonical forms on \(P\), compatibly with all transports.

Fix once and for all a K\"ahler form \(\omega_V\) on the ambient fourfold.
For a common resolution \(h:R\to P\), \(\rho:R\to Z\), define
\begin{equation}\label{eq:surface-class}
 c_Z=h_*[\rho^*(\omega_V|_Z)]\in H^{1,1}(P,\R).
\end{equation}
The brackets denote the de Rham class of the pulled-back local-potential
form on the smooth resolution.  All cohomology in the comparison below
is taken on smooth manifolds.
This is independent of further resolutions, since \(v_*v^*=1\) for a
modification \(v\).  Put \(\beta=[\rho^*(\omega_V|_Z)]\).  Its smooth
representative is semipositive and is strictly positive on a nonempty
open: the original stratum is generically immersed in the smooth
isomorphism locus of the special resolution.  Hence \(\beta^2>0\).
The blowup orthogonal decomposition has
\(\beta=h^*c_Z+e\), with \(e\) in the negative-definite exceptional
subspace.  Consequently
\begin{equation}\label{eq:surface-class-positive}
 c_Z^2=\beta^2-e^2>0.
\end{equation}

\begin{lemma}[The same ambient class across a link]\label{lem:class-link}
Let a link of a three-dimensional stratum compare non-Moishezon surface
strata \(Z_1,Z_2\), allowing \(Z_1=Z_2\).  Let \(P_i\) be their smooth
minimal models and \(\tau:P_1\simeq P_2\) the induced isomorphism.  For
the classes in \eqref{eq:surface-class},
\begin{equation}\label{eq:class-mod-ns}
 c_{Z_1}-\tau^*c_{Z_2}\in\NS(P_1)_\R.
\end{equation}
Here \(\NS(P)_\R\) is the real span of the first Chern classes of
holomorphic line bundles in \(H^{1,1}(P,\R)\).
\end{lemma}
\begin{proof}
Let \(i:Z\hookrightarrow V\) be the parent stratum, and choose a common
projective resolution with smooth compact K\"ahler source \(U\), with
maps \(a:U\to Z\) and \(b:U\to T\) to its original and Mori models.
Use the single class
\[
 \alpha=[(i\circ a)^*\omega_V]\in H^{1,1}(U,\R).
\]
Resolve the marked-branch comparison graph and its maps to \(U\) and
the minimal surfaces.  We obtain a smooth compact K\"ahler surface \(D\)
with maps \(r_i:D\to U\), bimeromorphic maps \(d_i:D\to P_i\), and
bimeromorphic maps to the original \(Z_i\), such that \(a r_i\) is the
inclusion of that original branch after its resolution and
\(d_2=\tau d_1\).  These identities hold on the
common marked open and hence everywhere.  This construction includes a
self-link: its two maps to the same original prime may differ by the
normalized exchange.  Resolution independence and projection formula give
\begin{equation}\label{eq:class-correspondence}
 c_{Z_1}-\tau^*c_{Z_2}
 =d_{1,*}(r_1^*-r_2^*)\alpha.
\end{equation}

The exceptional image of the modification \(b:U\to T\) has codimension
at least two in the normal threefold, and hence dimension at most one.
The Mori base in this two-marking case has dimension two.  After removing
its image, the discriminant, and the singular base locus, \(U\) is
therefore a smooth proper \(\PP^1\)-fibration over a dense open of the
base.  For such a fibration \(\pi:U^\circ\to W^\circ\), the differential
sequence and \(H^0(\PP^1,\Omega^1)=0\) give
\[
 \pi_*\Omega^2_{U^\circ}=\Omega^2_{W^\circ}.
\]
Indeed the quotient of \(\Omega^2_{U^\circ}\) by
\(\pi^*\Omega^2_{W^\circ}\) is
\(\pi^*\Omega^1_{W^\circ}\otimes\Omega^1_{U^\circ/W^\circ}\), whose
direct image is zero, while \(\pi_*\OO_{U^\circ}=\OO_{W^\circ}\).
Thus the two restrictions of a holomorphic two-form on \(U\) agree on
the dense paired-branch open of \(D\), even for the double branch after
an \'etale local ordering.  They agree on all of \(D\).

Holomorphic pullback and proper pushforward on compact K\"ahler manifolds
preserve Hodge type.  The map
\[
 d_{1,*}(r_1^*-r_2^*):H^2(U,\Q)\longrightarrow H^2(P_1,\Q)
\]
is therefore a rational Hodge morphism; pushforward here is between
equal-dimensional surfaces.  The preceding equality kills its
\((2,0)\) part and, by conjugation, its \((0,2)\) part.  Its rational
image is of type \((1,1)\), hence belongs to \(\NS(P_1)_\Q\) by
Lefschetz \((1,1)\).  Extend scalars to \(\R\) in
\eqref{eq:class-correspondence} to obtain \eqref{eq:class-mod-ns}.
\end{proof}

Swapnajit Das's positive-class and ruled-branch arguments
\cite[Lemmas~7.2--7.3 and Corollary~7.4]{Das2026} are close predecessors
of this mechanism for two disjoint degree-one branches.  The proof above
uses the same ambient class on both restrictions and includes the single
normalized degree-two branch.  It requires no rational polarization.

\subsection{A uniform scalar bound}

We spell out why the class congruence gives finite image for a whole group,
not just a finite-order scalar for each individually chosen comparison.

\begin{lemma}\label{lem:lattice-scalar}
Let \(P\) be a smooth connected nonprojective compact K\"ahler surface
with a nowhere-vanishing holomorphic two-form \(\eta\).  Put
\(N=\NS(P)_\R\).  Let \(G\subset\operatorname{Aut}(P)\) be a subgroup.
Assume either that \(N\) is degenerate, or that there exists
\(c\in H^{1,1}(P,\R)\) with \(c^2>0\) and
\(\sigma^*c-c\in N\) for every \(\sigma\in G\).  Then the image of
\(G\) on \(\C\eta^{\otimes m}\) is finite for every \(m>0\).
More precisely, every volume scalar has order in the finite set of
integers \(n\) with \(\varphi(n)\leq b_2(P)\).
\end{lemma}
\begin{proof}
The intersection form on \(H^{1,1}(P,\R)\) has Lorentz signature
\((1,h^{1,1}-1)\).  Its restriction to \(N\) is nonpositive.  To see
this, a rational positive-square class in \(N\), after scaling and
choosing its sign, is \(c_1(L)\) with positive K\"ahler degree.
Riemann--Roch and Serre duality give \(h^0(P,L^v)\gg v^2\): the Euler
characteristic has positive quadratic leading term and
\(H^0(P,K_P\otimes L^{-v})=0\) for large \(v\) by its negative K\"ahler
degree.  This would make \(P\) Moishezon and hence projective.  Rational
approximation in \(N\) rules out a positive real class as well.

If \(N\) is negative definite, orthogonally project \(c\) to \(N^\perp\).
The projection \(c_\perp\) has positive square and is fixed by every
\(\sigma^*\): the congruence kills its projected difference, and each
automorphism preserves \(N\) and the intersection form.  The perpendicular
of \(c_\perp\) is negative definite, so all eigenvalues on
\(H^{1,1}(P,\R)\) have modulus one.

If \(N\) is degenerate, its radical is a rational isotropic line
\(\ell\); nonpositivity in a Lorentz space permits no larger radical.
An integral automorphism preserves a primitive integral generator of
\(\ell\) up to sign.  The invariant flag
\[
 \ell\subset\ell^\perp\subset H^{1,1}(P,\R)
\]
has eigenvalues \(\pm1\) on the first and last quotients, which pair
dually, and a negative-definite middle quotient \(\ell^\perp/\ell\).
All \((1,1)\) eigenvalues again have modulus one.  This argument permits
unipotent action and makes no finiteness assertion for all cohomology.

Write \(\sigma^*\eta=\delta_\sigma\eta\).  Integration of
\(\eta\wedge\bar\eta\) gives \(|\delta_\sigma|=1\), and the conjugate
eigenvalue has the same modulus.  Since a nowhere-vanishing canonical form
spans \(H^{2,0}(P)\), all eigenvalues on \(H^2(P,\C)\) now have modulus
one.  The action on \(H^2(P,\Z)/\mathrm{torsion}\) is integral.  Kronecker's
theorem makes each eigenvalue a root of unity; if its order is \(n\), its
cyclotomic polynomial has degree \(\varphi(n)\leq b_2(P)\).  There are
only finitely many such \(n\).  The scalars of every element of \(G\)
therefore lie in one finite set of roots of unity, which proves finite
image on \(\C\eta^{\otimes m}\).
\end{proof}

If \(P\) has algebraic dimension zero, its minimality and
\cite[Theorem~4]{KodairaCompactAnalytic} imply that its canonical line
is trivial, so it has a nowhere-vanishing
volume \(\eta\).  Every meromorphic pluriform is
a constant multiple of a power of \(\eta\), since the quotient is a
meromorphic function and \(a(P)=0\).  In particular the section space
under consideration has dimension at most one.  Composing
\eqref{eq:class-mod-ns} along a returning sequence of links gives
\(\sigma^*c_Z-c_Z\in\NS(P)_\R\); each intervening isomorphism preserves
N\'eron--Severi.  Lemma~\ref{lem:lattice-scalar} applies, with the fixed
cohomology rank of this \(P\).  It gives a finite scalar image for all
returning compositions at once.

\subsection{Algebraic dimension one}

Suppose \(a(P)=1\).  Its holomorphic algebraic reduction is an elliptic
fibration \(v:P\to J\) with connected fibers over a smooth compact curve,
and every curve on \(P\) is vertical
\cite[Theorem~2]{KodairaCompactAnalytic}.  Every automorphism preserves
this reduction, which is determined by the meromorphic function field.
Let \(G\) be the group of returning link automorphisms and let \(H\) be
the corresponding finite-dimensional invariant space of meromorphic
\(qk\)-pluriforms on \(P\).

On a smooth base open avoiding the finitely many polar fibers of a basis
of \(H\), division by a base differential makes these forms holomorphic
powers of the elliptic differential on each fiber.  An automorphism over
the base multiplies that differential by a root of unity of order
\(1,2,3,4\), or \(6\); its multiplier is holomorphic with values in a
finite set, hence locally constant.  It acts on all of \(H\) by the
corresponding common scalar.  Thus the kernel of \(G\to\operatorname{Aut}(J)\)
has finite image on \(H\).  If the image on \(J\) is finite, its finitely
many cosets make the full image finite as well.

It remains to consider an infinite base image.  A curve of genus at
least two has finite automorphism group.  If \(J=\PP^1\), the finite set
of nonsmooth fibers, including multiple fibers, is invariant, so it would
have at most two points.  This contradicts the theorem that a nonalgebraic compact
K\"ahler elliptic surface over \(\PP^1\) has at least three singular
fibers \cite[Proposition~A.1]{ClaudonHoeringLin2019}.  This theorem is a
short form of the elliptic canonical-degree exclusion needed here and
does not assume a section of the fibration.

If \(J\) has genus one, the infinite base group contains infinitely many
translations.  Its invariant finite special-fiber set must be empty.
Choose a nonzero holomorphic two-form \(\eta\) on \(P\), whose existence
was proved above, and write its effective integral zero divisor as
\(D_P\).  Every component is vertical.  Each fiber is now smooth and
connected, hence an irreducible reduced elliptic curve; a vertical prime
is that entire fiber, and smoothness gives it multiplicity one in the
pullback of its base point.  Thus \(D_P=v^*D_J\) for an effective integral
divisor \(D_J\) on \(J\), and the actual section gives
\[
 \omega_P\simeq\OO_P(D_P)\simeq v^*\OO_J(D_J).
\]
Pullback \(v^*:\operatorname{Pic}(J)\to\operatorname{Pic}(P)\) is
injective.  Indeed, if \(v^*M\simeq\OO_P\), connected fibers and
projection formula give
\[
 M\simeq M\otimes v_*\OO_P\simeq v_*v^*M\simeq v_*\OO_P\simeq\OO_J.
\]
For an automorphism \(\sigma\in G\) covering \(h\in\operatorname{Aut}(J)\),
the identities \(v\sigma=hv\) and \(\sigma^*\omega_P\simeq\omega_P\)
therefore imply \(h^*\OO_J(D_J)\simeq\OO_J(D_J)\).

Put \(d=\deg D_J\) and write \(J=\C/\Lambda\).  If translation \(t_x\)
stabilizes \(\OO_J(D_J)\), then \(t_x^*D_J-D_J\) is principal.  Its
point sum in \(J\) is \(-dx\), whereas a principal divisor has point
sum zero.  To recall the latter fact, lift its meromorphic function to
an elliptic function \(F\), and choose a fundamental parallelogram
\(\mathcal P\) with boundary avoiding its zeros and poles.  For a lattice
basis \(\omega_1,\omega_2\), the residue theorem and pairing opposite
edges give
\[
 \sum_{z\in\mathcal P}\ord_z(F)z
 =\frac{1}{2\pi i}\int_{\partial\mathcal P}z\frac{F'(z)}{F(z)}\,dz
 =\omega_1 n_2-\omega_2 n_1\in\Lambda,
\]
where \(n_i=(2\pi i)^{-1}\int_{z_0}^{z_0+\omega_i}F'(z)/F(z)\,dz\in\Z\), since
the endpoint values of \(F\) agree.  It follows that \(dx=0\) in \(J\).
If \(d>0\), all stabilizing translations lie in the finite group
\(J[d]=(\tfrac1d\Lambda)/\Lambda\).  The infinitely many translations
above therefore force \(d=0\).  Effectivity gives \(D_J=0\), so \(\eta\)
is nowhere vanishing.

The ratio of any holomorphic two-form to \(\eta\) is holomorphic on the
compact connected \(P\), hence constant.  Thus \(\eta\) spans
\(H^0(P,\omega_P)\), and every \(\sigma\in G\) satisfies
\(\sigma^*\eta=c_\sigma\eta\) for a nonzero constant \(c_\sigma\).
Every element of \(H\) is \(\eta^{\otimes qk}\)
times a meromorphic function from \(J\).  The corresponding function
space is therefore preserved by the base action.  The union of the pole sets of
a basis is finite and intrinsic to this vector space, hence invariant
under the base group.  Infinitely many translations preserve no nonempty
finite subset.  Thus these functions have no poles and are constant.

The fiber class belongs to \(\NS(P)_\R\), has square zero, and is nonzero
by its positive K\"ahler area.  Nonpositivity of \(\NS(P)_\R\) makes it
a radical vector.  The degenerate case of
Lemma~\ref{lem:lattice-scalar} gives the same finite scalar image on \(H\).
This completes the algebraic-dimension-one case and the proof of
Proposition~\ref{prop:finite-isotropy}.

For later use, finiteness of isotropy also controls all transports in a
fixed degree between two objects of the same orbit.  Choose one comparison
from an orbit representative to each member.  Every other transport to
that member is a self-transport of the representative, whose image is
finite, followed by this chosen transport.  Thus only finitely many linear
transport actions occur in that degree, even if the groupoid has infinitely
many bimeromorphic arrows.
